\documentclass[border=5mm]{standalone}
% ==============================================================================
% SETUP.TEX - CONFIGURACIÓN GLOBAL DEL PROYECTO
% ==============================================================================
% Este archivo centraliza todos los paquetes y estilos.
% Debe incluirse en el preámbulo del libro principal y de los archivos standalone.
% \PassOptionsToPackage{gray}{xcolor}
% --- 1. CODIFICACIÓN E IDIOMA ---
% fix-cm habilita las fuentes Computer Modern en tamaños arbitrarios (por debajo
% de 5 pt, entre otros). Debe cargarse antes que fontenc.
\usepackage{fix-cm}
\usepackage[utf8]{inputenc}
\usepackage[T1]{fontenc}
\usepackage[spanish, es-tabla]{babel}  
\usepackage{csquotes} % Recomendado para babel y citas

\usepackage{import}
% mode=image: Usa el PDF pre-compilado, nunca intenta recompilar.
% Debes compilar cada figura manualmente antes de compilar el libro.
\usepackage[mode=image, subpreambles=false]{standalone}

% --- 2. MATEMÁTICAS Y CIENCIAS ---
\usepackage{amsmath, amssymb, amsfonts, mathtools}
\usepackage{enumitem}

% LaTeX trae \DeclareMathSizes{5}{5}{5}{5}, así que a 5 pt (\tiny) los subíndices
% salen del mismo tamaño que la base: en las etiquetas de figuras el M de
% $\omega_M$ queda desproporcionado. Se restablece la proporción habitual.
\DeclareMathSizes{5}{5}{3.7}{3}

% --- 3. DISEÑO Y GEOMETRÍA --- 
% Unificamos los márgenes para todo el documento
\usepackage[left=2.5cm,right=2.5cm,top=3cm,bottom=3cm]{geometry}
\makeatletter
\ifstandalone % Evita que geometry dañe el recorte de las figuras bajándolas 3cm
  \@ifclasswith{standalone}{crop=false}{
    % Es un capítulo/sección: Dejamos que geometry actúe.
  }{
    % Es una figura: Neutralizamos geometry para que no las recorte mal.
    \geometry{pass}
  }
\fi
\makeatother
\usepackage{parskip} % Espaciado entre párrafos sin sangría
\usepackage{float}   % Para usar [H] en figura

% Ajustes de flotantes para evitar espacios verticales exagerados
\renewcommand{\topfraction}{0.95}      % Permite que las figuras ocupen hasta el 95% superior
\renewcommand{\bottomfraction}{0.95}   % Permite que las figuras ocupen hasta el 95% inferior
\renewcommand{\textfraction}{0.05}     % Permite hojas con tan solo un 5% de texto
\renewcommand{\floatpagefraction}{0.8} % Requiere que una página de solo figuras esté 80% llena
\setcounter{topnumber}{4}
\setcounter{bottomnumber}{4}
\setcounter{totalnumber}{10}

% --- 4. GRÁFICOS Y TABLAS ---
\usepackage{graphicx}
\usepackage{booktabs} % Tablas profesionales
\usepackage{caption}  % Control de captions y \captionof
\usepackage{tikz}
\usetikzlibrary{arrows.meta, calc, angles, quotes, babel, backgrounds}
\usepackage{pgfplots}
\usepgfplotslibrary{groupplots}
\usepgfplotslibrary{external}
\usepgfplotslibrary{patchplots}
\pgfplotsset{compat=1.18}
\usepackage[american, straightvoltages]{circuitikz}

% --- 5. COLORES Y CAJAS ---

\usepackage[table]{xcolor}
\usepackage{tcolorbox}

% Definición de colores corporativos/estilo
\definecolor{utnblue}{RGB}{0, 51, 102}
\definecolor{codegreen}{rgb}{0,0.6,0}
\definecolor{codegray}{rgb}{0.5,0.5,0.5}
\definecolor{codepurple}{rgb}{0.58,0,0.82}
\definecolor{backcolour}{rgb}{0.96,0.96,0.96}

% --- 6. ENTORNOS PERSONALIZADOS (CAJAS) ---

% Caja "Resultado" (Azul Corporativo) — Definiciones, fórmulas clave, resultados importantes.
% Uso: \begin{resultado}[Fórmula de Euler] ... \end{resultado}
\newtcolorbox{resultado}[1][]{
    colback=utnblue!5!white,
    colframe=utnblue,
    title=\textbf{#1},
    fonttitle=\bfseries,
    sharp corners,
    boxrule=0.5mm
}

% Caja "Nota" (Gris) — Advertencias, aplicaciones, contexto secundario.
% Uso: \begin{nota}[Cuidado con la calculadora] ... \end{nota}
% Uso: \begin{nota} ... \end{nota}  → título por defecto "Nota"
\newtcolorbox{nota}[1][Nota]{
    colback=codegray!5!white,
    colframe=codegray,
    title=\textbf{#1},
    fonttitle=\bfseries,
    sharp corners,
    boxrule=0.5mm
}

% Caja inline para resaltar un valor y enlazarlo con su otra aparición en una tabla
% o en una cuenta: mismo color, mismo valor.
% Uso: \marcavalor{$4$}  o  \marcavalor[codegreen]{$4$}
\newtcbox{\marcavalor}[1][utnblue]{
    on line,
    boxsep=1pt, left=2pt, right=2pt, top=1pt, bottom=1pt,
    colback=#1!10!white,
    colframe=#1,
    boxrule=0.4pt,
    arc=2pt
}

% --- 7. CÓDIGO FUENTE (LISTINGS) ---
\usepackage{listings}

\lstdefinestyle{miestilo}{
    backgroundcolor=\color{backcolour},   
    commentstyle=\color{codegreen},
    keywordstyle=\color{magenta},
    numberstyle=\tiny\color{codegray},
    stringstyle=\color{codepurple},
    basicstyle=\ttfamily\footnotesize,
    breakatwhitespace=false,         
    breaklines=true,                 
    captionpos=b,                    
    keepspaces=true,                 
    numbers=left,                    
    numbersep=5pt,                  
    showspaces=false,                
    showstringspaces=false,
    showtabs=false,                  
    tabsize=2,
    frame=single,                    
    rulecolor=\color{codegray}
}

\lstset{style=miestilo}
% Soporte para tildes y ñ en bloques de código
\lstset{literate={ñ}{{\~n}}1 {á}{{\'a}}1 {é}{{\'e}}1 {í}{{\'i}}1 {ó}{{\'o}}1 {ú}{{\'u}}1}

% Operadores matemáticos personalizados

% Vector en sentido discreto (lista finita de coordenadas): negrita + flecha.
% Uso: \vect{v}, \vect{e}_k, \vect{0}.
\newcommand{\vect}[1]{\vec{\mathbf{#1}}}

% Fecha de versión y comando para capítulos standalone
\providecommand{\verdate}{\the\year.\ifnum\month<10 0\fi\the\month.\ifnum\day<10 0\fi\the\day}
\newcommand{\chapterversion}{%
    \vspace{-1.5cm}\begin{flushright}\small\textit{\color{codegray}Versión~\verdate}\end{flushright}\vspace{0.5cm}%
}

% --- 8. HIPERVÍNCULOS (DEBE SER EL ÚLTIMO PAQUETE) ---
\usepackage[colorlinks=true, linkcolor=utnblue, urlcolor=utnblue, citecolor=utnblue]{hyperref}

% --- 9. REFERENCIAS CRUZADAS ENTRE CAPÍTULOS STANDALONE ---
% Cuando se compila un capítulo standalone, xr-hyper lee los .aux de los
% otros capítulos (si existen) para resolver \ref{} a labels externos.
% En la compilación del libro completo este bloque no se activa.
\ifstandalone
  \usepackage{xr-hyper}
  \makeatletter
  \newcommand{\xrchap}[1]{%
    \edef\xrc@self{\detokenize{Capitulo#1}}%
    \edef\xrc@job{\jobname}%
    \ifx\xrc@self\xrc@job\else
      \IfFileExists{../#1capitulo/Capitulo#1.aux}%
        {\externaldocument{../#1capitulo/Capitulo#1}}{}%
    \fi
  }
  \makeatother
  \xrchap{01}\xrchap{02}\xrchap{03}\xrchap{04}
  \xrchap{05}\xrchap{06}\xrchap{07}\xrchap{08}
  \xrchap{09}\xrchap{10}\xrchap{11}\xrchap{12}
  \xrchap{13}
\fi
\usepgfplotslibrary{fillbetween}
% \PassOptionsToPackage{gray}{xcolor}

% Superposición de dos respuestas al impulso.
% Parámetros: h(t) = e^{-t} u(t), t1 = 0, alpha1 = 0.5, t2 = 1, alpha2 = 1.
%
% Layout 3x2:
%   (a)  x1 = alpha1 δ(t-t1)         (a') y1 = alpha1 h(t-t1)
%   (b)  x2 = alpha2 δ(t-t2)         (b') y2 = alpha2 h(t-t2)
%   (c)  x = x1 + x2                 (c') y = y1 + y2  (aportes apilados)

\begin{document}
\begin{tikzpicture}[>=Latex, font=\small]

\colorlet{col1}{utnblue}
\colorlet{col2}{red!80!black}
\colorlet{col1light}{utnblue!25}
\colorlet{col2light}{red!80!black!30}
\colorlet{tlabel}{gray!50!black}

\pgfplotsset{
    every axis/.append style={clip=false},
}

% =============== Fila 1: x1 — y1 ===============

\begin{axis}[
    name=a1,
    width=7cm, height=3.2cm,
    axis lines=center,
    xmin=-0.4, xmax=5.3,
    ymin=-0.25, ymax=1.25,
    xtick={1}, xticklabels={$t_2$},
    ytick=\empty,
    tick style={thin},
    xlabel={$t$}, xlabel style={anchor=north west, at={(axis description cs:1,0)}},
    ylabel={$x_1$},
    ylabel style={font=\small, rotate=0, anchor=south east,
                  at={(axis description cs:0,1)}},
    title={\footnotesize (a) $x_1(t) = \alpha_1\,\delta(t-t_1)$},
    title style={font=\footnotesize\bfseries},
    every tick label/.style={font=\scriptsize},
]
    \draw[->, col1, very thick, line width=1.2pt]
        (axis cs:0,0) -- (axis cs:0,0.5);
    \node[col1, font=\scriptsize, anchor=west] at (axis cs:0.12, 0.42) {$\alpha_1$};
    \node[font=\scriptsize, anchor=base, tlabel, xshift=5pt, yshift=-10pt] at (axis cs:0,0) {$t_1$};
    \draw[dashed, gray!55, thin] (axis cs:1, 0) -- (axis cs:1, 1.15);
\end{axis}

\begin{axis}[
    name=b1,
    at={(a1.east)}, anchor=west, xshift=1.1cm,
    width=7cm, height=3.2cm,
    axis lines=center,
    xmin=-0.4, xmax=5.3,
    ymin=-0.25, ymax=1.25,
    xtick={1}, xticklabels={$t_2$},
    ytick=\empty,
    tick style={thin},
    xlabel={$t$}, xlabel style={anchor=north west, at={(axis description cs:1,0)}},
    ylabel={$y_1$},
    ylabel style={font=\small, rotate=0, anchor=south east,
                  at={(axis description cs:0,1)}},
    title={\footnotesize (a$'$) $y_1(t) = \alpha_1\,h(t-t_1)$},
    title style={font=\footnotesize\bfseries},
    every tick label/.style={font=\scriptsize},
]
    \addplot[col1light, draw=none, domain=0:5.2, samples=120] {0.5*exp(-x)} \closedcycle;
    \addplot[col1, very thick, domain=-0.4:0, samples=2] {0};
    \addplot[col1, very thick, domain=0:5.2, samples=120] {0.5*exp(-x)};
    \draw[col1, very thick] (axis cs:0,0) -- (axis cs:0,0.5);
    \fill[col1]  (axis cs:0,0.5) circle (1.8pt);
    \fill[white] (axis cs:0,0) circle (1.8pt);
    \draw[col1]  (axis cs:0,0) circle (1.8pt);
    % Línea horizontal marcando la amplitud inicial alpha_1
    \draw[col1, thin] (axis cs:-0.4, 0.5) -- (axis cs:5.3, 0.5);
    \node[col1, font=\scriptsize, anchor=south east] at (axis cs:5.25, 0.5) {$\alpha_1$};
    \node[font=\scriptsize, anchor=base, tlabel, xshift=5pt, yshift=-10pt] at (axis cs:0,0) {$t_1$};
    \draw[dashed, gray!55, thin] (axis cs:1, 0) -- (axis cs:1, 1.15);
\end{axis}

% =============== Fila 2: x2 — y2 ===============

\begin{axis}[
    name=a2,
    at={(a1.south)}, anchor=north, yshift=-1.2cm,
    width=7cm, height=3.2cm,
    axis lines=center,
    xmin=-0.4, xmax=5.3,
    ymin=-0.25, ymax=1.25,
    xtick={1}, xticklabels={$t_2$},
    ytick=\empty,
    tick style={thin},
    xlabel={$t$}, xlabel style={anchor=north west, at={(axis description cs:1,0)}},
    ylabel={$x_2$},
    ylabel style={font=\small, rotate=0, anchor=south east,
                  at={(axis description cs:0,1)}},
    title={\footnotesize (b) $x_2(t) = \alpha_2\,\delta(t-t_2)$},
    title style={font=\footnotesize\bfseries},
    every tick label/.style={font=\scriptsize},
]
    \draw[->, col2, very thick, line width=1.2pt]
        (axis cs:1,0) -- (axis cs:1,1);
    \node[col2, font=\scriptsize, anchor=west] at (axis cs:1.62, 0.88) {$\alpha_2$};
    \node[font=\scriptsize, anchor=base, tlabel, xshift=5pt, yshift=-10pt] at (axis cs:0,0) {$t_1$};
    \draw[dashed, gray!55, thin] (axis cs:0, 0) -- (axis cs:0, 1.15);
\end{axis}

\begin{axis}[
    name=b2,
    at={(b1.south)}, anchor=north, yshift=-1.2cm,
    width=7cm, height=3.2cm,
    axis lines=center,
    xmin=-0.4, xmax=5.3,
    ymin=-0.25, ymax=1.25,
    xtick={1}, xticklabels={$t_2$},
    ytick=\empty,
    tick style={thin},
    xlabel={$t$}, xlabel style={anchor=north west, at={(axis description cs:1,0)}},
    ylabel={$y_2$}, 
    ylabel style={font=\small, rotate=0, anchor=south east,
                  at={(axis description cs:0,1)}},
    title={\footnotesize (b$'$) $y_2(t) = \alpha_2\,h(t-t_2)$},
    title style={font=\footnotesize\bfseries},
    every tick label/.style={font=\scriptsize},
]
    \addplot[col2light, draw=none, domain=1:5.2, samples=120] {exp(-(x-1))} \closedcycle;
    \addplot[col2, very thick, domain=-0.4:1, samples=2] {0};
    \addplot[col2, very thick, domain=1:5.2, samples=120] {exp(-(x-1))};
    \draw[col2, very thick] (axis cs:1,0) -- (axis cs:1,1);
    \fill[col2]  (axis cs:1,1) circle (1.8pt);
    \fill[white] (axis cs:1,0) circle (1.8pt);
    \draw[col2]  (axis cs:1,0) circle (1.8pt);
    % Línea horizontal marcando la amplitud inicial alpha_2
    \draw[col2, thin] (axis cs:-0.4, 1) -- (axis cs:5.3, 1);
    \node[col2, font=\scriptsize, anchor=south east] at (axis cs:5.25, 1) {$\alpha_2$};
    \node[font=\scriptsize, anchor=base, tlabel, xshift=5pt, yshift=-10pt] at (axis cs:0,0) {$t_1$};
    \draw[dashed, gray!55, thin] (axis cs:0, 0) -- (axis cs:0, 1.15);
\end{axis}

% =============== Fila 3: suma entrada — suma salida con aportes apilados ===============

\begin{axis}[
    name=a3,
    at={(a2.south)}, anchor=north, yshift=-1.2cm,
    width=7cm, height=3.2cm,
    axis lines=center,
    xmin=-0.4, xmax=5.3,
    ymin=-0.25, ymax=1.25,
    xtick={1}, xticklabels={$t_2$},
    ytick=\empty,
    tick style={thin},
    xlabel={$t$}, xlabel style={anchor=north west, at={(axis description cs:1,0)}},
    ylabel={$x$},
    ylabel style={font=\small, rotate=0, anchor=south east,
                  at={(axis description cs:0,1)}},
    title={\footnotesize (c) $x(t) = x_1(t) + x_2(t)$},
    title style={font=\footnotesize\bfseries},
    every tick label/.style={font=\scriptsize},
]
    \draw[->, col1, very thick, line width=1.2pt]
        (axis cs:0,0) -- (axis cs:0,0.5);
    \node[col1, font=\scriptsize, anchor=west] at (axis cs:0.12, 0.42) {$\alpha_1$};
    \draw[->, col2, very thick, line width=1.2pt]
        (axis cs:1,0) -- (axis cs:1,1);
    \node[col2, font=\scriptsize, anchor=west] at (axis cs:1.62, 0.88) {$\alpha_2$};
    \node[font=\scriptsize, anchor=base, tlabel, xshift=5pt, yshift=-10pt] at (axis cs:0,0) {$t_1$};
\end{axis}

\begin{axis}[
    name=b3,
    at={(b2.south)}, anchor=north, yshift=-1.2cm,
    width=7cm, height=3.2cm,
    axis lines=center,
    xmin=-0.4, xmax=5.3,
    ymin=-0.25, ymax=1.25,
    xtick={1}, xticklabels={$t_2$},
    ytick=\empty,
    tick style={thin},
    xlabel={$t$}, xlabel style={anchor=north west, at={(axis description cs:1,0)}},
    ylabel={$y$},
    ylabel style={font=\small, rotate=0, anchor=south east,
                  at={(axis description cs:0,1)}},
    title={\footnotesize (c$'$) $y(t) = y_1(t) + y_2(t)$},
    title style={font=\footnotesize\bfseries},
    every tick label/.style={font=\scriptsize},
]
    % ---------- Paso 1: calcular y sombrear el área bajo y_total = y1 + y2 ----------
    % Para t < t2, y_total = y1 (rojo y azul coinciden en esa zona).
    \addplot[fill=col2light, draw=none, domain=0:1, samples=80] {0.5*exp(-x)} \closedcycle;
    % Para t >= t2, y_total = y1 + y2.
    \addplot[fill=col2light, draw=none, domain=1:5.2, samples=120]
        {0.5*exp(-x) + exp(-(x-1))} \closedcycle;

    % ---------- Paso 2: sombrear encima el área bajo y1 ----------
    % Ocluye al rojo donde y1 > 0; deja visible sólo la banda y_total \ y1.
    \addplot[fill=col1light, draw=none, domain=0:5.2, samples=120] {0.5*exp(-x)} \closedcycle;

    % ---------- Envelope y_total = y1 + y2 (negro, grueso) ----------
    \addplot[black, very thick, domain=-0.4:0, samples=2] {0};
    \addplot[black, very thick, domain=0:1, samples=80] {0.5*exp(-x)};
    \addplot[black, very thick, domain=1:5.2, samples=160]
        {0.5*exp(-x) + exp(-(x-1))};

    % Discontinuidad en t=t1=0
    \draw[black, very thick] (axis cs:0,0) -- (axis cs:0,0.5);
    \fill[black]  (axis cs:0,0.5) circle (1.8pt);
    \fill[white]  (axis cs:0,0) circle (1.8pt);
    \draw[black]  (axis cs:0,0) circle (1.8pt);

    % Discontinuidad en t=t2=1 (salto = alpha_2 = 1)
    \pgfmathsetmacro{\yleft}{0.5*exp(-1)}
    \pgfmathsetmacro{\yright}{0.5*exp(-1) + 1}
    \draw[black, very thick] (axis cs:1,\yleft) -- (axis cs:1,\yright);
    \fill[black]  (axis cs:1,\yright) circle (1.8pt);
    \fill[white]  (axis cs:1,\yleft) circle (1.8pt);
    \draw[black]  (axis cs:1,\yleft) circle (1.8pt);

    % Líneas horizontales de referencia (amplitudes iniciales), dashed en este panel
    \draw[col1, thin, dashed] (axis cs:-0.4, 0.5) -- (axis cs:5.3, 0.5);
    \node[col1, font=\scriptsize, anchor=south east] at (axis cs:5.25, 0.5) {$\alpha_1$};
    \draw[col2, thin, dashed] (axis cs:-0.4, 1) -- (axis cs:5.3, 1);
    \node[col2, font=\scriptsize, anchor=south east] at (axis cs:5.25, 1) {$\alpha_2$};

    % Etiqueta del envelope
    \node[font=\scriptsize, anchor=south west] at (axis cs:2.6, 0.70) {$y = y_1 + y_2$};

    \node[font=\scriptsize, anchor=base, tlabel, xshift=5pt, yshift=-10pt] at (axis cs:0,0) {$t_1$};
\end{axis}

\end{tikzpicture}
\end{document}
